% GENERATED FILE. Edit canonical lessons, then run npm run generate:books.
% canonical-source-hash: fnv1a64:80b53788251d6d65
% canonical-lessons: SA-C426-videsah, SA-C426-janakirnah, SA-C426-sulkasthanam, SA-C426-avatisthate, SA-C426-pratisrayah

\chapter{Abroad}
\label{ch:sa-abroad}
\hlchaptermodality{426}

\begin{chapteropening}
\textbf{By the end of this chapter:} \emph{I can talk about being abroad, the crowds and the customs house.}

Complete the last lesson of chapter 426: I can talk about being abroad, the crowds and the customs house.
\end{chapteropening}

\section[\texorpdfstring{\sk{विदेशः} (videśaḥ)}{विदेशः (videśaḥ)}]{\sk{विदेशः} (videśaḥ) — a foreign country; abroad}

\label{lesson:SA-C426-videsah}



\begin{quote}
Before the new one: say the Sanskrit for a round window, then the Sanskrit for memorable.
\end{quote}

\subsection*{You'll want to know: \sk{विदेशः}}
\textbf{\sk{विदेशः}} — \emph{videśaḥ} — \textquotedblleft{}a foreign country; abroad\textquotedblright{}.

\begin{quote}
\sk{मम} \sk{भ्राता} \sk{विदेशे} \sk{वसति।}
\end{quote}

\begin{quote}
\emph{mama bhrātā videśe vasati.}
\end{quote}

\begin{quote}
My brother lives abroad.
\end{quote}

A masculine noun.

\begin{cousinweb}
\textbf{\sk{विदेशः}} (\emph{videśaḥ}) is built from \emph{vi}, \textquotedblleft{}apart, away\textquotedblright{}, and \emph{deśa}, \textquotedblleft{}a country\textquotedblright{}; MW gives \emph{videśa} as another country, a foreign land, abroad (MW).
\end{cousinweb}

\begin{grammarlens}[title={What you've built}]
One more word for getting around.
\end{grammarlens}

\subsection*{Your turn}
\begin{itemize}
\raggedright
  \item \emph{Say it:} \emph{videśaḥ}
  \item \emph{Say it:} \emph{videśaḥ}, once more
  \item \emph{Say it:} \emph{mama bhrātā videśe vasati.} — My brother lives abroad.
  \item \emph{Recall:} say the Sanskrit for a round window, then the Sanskrit for memorable, then say \emph{videśaḥ} again
\end{itemize}

\begin{tcolorbox}[breakable,colback=teal!4,colframe=teal!35!black,title={Before you move on}]
What does \sk{विदेशः} mean? (\textquotedblleft{}A foreign country; abroad\textquotedblright{}.) Say it once more.
\end{tcolorbox}
\section[\texorpdfstring{\sk{जनाकीर्णः} (janākīrṇaḥ)}{जनाकीर्णः (janākīrṇaḥ)}]{\sk{जनाकीर्णः} (janākīrṇaḥ) — crowded}

\label{lesson:SA-C426-janakirnah}



\begin{quote}
Before the new one: say the Sanskrit for memorable, then the Sanskrit for a foreign country.
\end{quote}

\subsection*{You'll want to know: \sk{जनाकीर्णः}}
\textbf{\sk{जनाकीर्णः}} — \emph{janākīrṇaḥ} — \textquotedblleft{}crowded\textquotedblright{}.

\begin{quote}
\sk{विदेशस्य} \sk{नगराणि} \sk{जनाकीर्णानि} \sk{सन्ति।}
\end{quote}

\begin{quote}
\emph{videśasya nagarāṇi janākīrṇāni santi.}
\end{quote}

\begin{quote}
The cities abroad are crowded.
\end{quote}

\begin{cousinweb}
\textbf{\sk{जनाकीर्णः}} (\emph{janākīrṇaḥ}) is built from \emph{jana}, \textquotedblleft{}people\textquotedblright{}, and \emph{ākīrṇa}, \textquotedblleft{}scattered over, filled\textquotedblright{}, from \emph{ā} and the root √kṝ, \textquotedblleft{}to scatter\textquotedblright{}; MW gives the compound as crowded with people (MW).
\end{cousinweb}

\begin{grammarlens}[title={What you've built}]
One more word for getting around.
\end{grammarlens}

\subsection*{Your turn}
\begin{itemize}
\raggedright
  \item \emph{Say it:} \emph{janākīrṇaḥ}
  \item \emph{Say it:} \emph{janākīrṇaḥ}, once more
  \item \emph{Say it:} \emph{videśasya nagarāṇi janākīrṇāni santi.} — The cities abroad are crowded.
  \item \emph{Recall:} say the Sanskrit for memorable, then the Sanskrit for a foreign country, then say \emph{janākīrṇaḥ} again
\end{itemize}

\begin{tcolorbox}[breakable,colback=teal!4,colframe=teal!35!black,title={Before you move on}]
What does \sk{जनाकीर्णः} mean? (\textquotedblleft{}Crowded\textquotedblright{}.) Say it once more.
\end{tcolorbox}
\section[\texorpdfstring{\sk{शुल्कस्थानम्} (śulkasthānam)}{शुल्कस्थानम् (śulkasthānam)}]{\sk{शुल्कस्थानम्} (śulkasthānam) — a customs house}

\label{lesson:SA-C426-sulkasthanam}



\begin{quote}
Before the new one: say the Sanskrit for a foreign country, then the Sanskrit for crowded.
\end{quote}

\subsection*{You'll want to know: \sk{शुल्कस्थानम्}}
\textbf{\sk{शुल्कस्थानम्}} — \emph{śulkasthānam} — \textquotedblleft{}a customs house\textquotedblright{}.

\begin{quote}
\sk{सीमान्ते} \sk{शुल्कस्थानम्} \sk{अस्ति।}
\end{quote}

\begin{quote}
\emph{sīmānte śulkasthānam asti.}
\end{quote}

\begin{quote}
There is a customs house at the border.
\end{quote}

A neuter noun.

\begin{cousinweb}
\textbf{\sk{शुल्कस्थानम्}} (\emph{śulkasthānam}) is built from \emph{śulka}, \textquotedblleft{}a toll, a duty\textquotedblright{}, and \emph{sthāna}, \textquotedblleft{}a place\textquotedblright{}; MW gives the compound as a toll-house, a tax office, a custom house (MW).
\end{cousinweb}

\begin{grammarlens}[title={What you've built}]
One more word for getting around.
\end{grammarlens}

\subsection*{Your turn}
\begin{itemize}
\raggedright
  \item \emph{Say it:} \emph{śulkasthānam}
  \item \emph{Say it:} \emph{śulkasthānam}, once more
  \item \emph{Say it:} \emph{sīmānte śulkasthānam asti.} — There is a customs house at the border.
  \item \emph{Recall:} say the Sanskrit for a foreign country, then the Sanskrit for crowded, then say \emph{śulkasthānam} again
\end{itemize}

\begin{tcolorbox}[breakable,colback=teal!4,colframe=teal!35!black,title={Before you move on}]
What does \sk{शुल्कस्थानम्} mean? (\textquotedblleft{}A customs house\textquotedblright{}.) Say it once more.
\end{tcolorbox}
\section[\texorpdfstring{\sk{अवतिष्ठते} (avatiṣṭhate)}{अवतिष्ठते (avatiṣṭhate)}]{\sk{अवतिष्ठते} (avatiṣṭhate) — stays on; remains}

\label{lesson:SA-C426-avatisthate}



\begin{quote}
Before the new one: say the Sanskrit for crowded, then the Sanskrit for a customs house.
\end{quote}

\subsection*{You'll want to know: \sk{अवतिष्ठते}}
\textbf{\sk{अवतिष्ठते}} — \emph{avatiṣṭhate} — \textquotedblleft{}stays on; remains\textquotedblright{}.

\begin{quote}
\sk{सः} \sk{विदेशे} \sk{एकं} \sk{संवत्सरम्} \sk{अवतिष्ठते।}
\end{quote}

\begin{quote}
\emph{saḥ videśe ekaṃ saṃvatsaram avatiṣṭhate.}
\end{quote}

\begin{quote}
He stays on abroad for a year.
\end{quote}

Past (he …): \sk{सः} \sk{विदेशे} \sk{अवस्थितवान्}, \emph{saḥ videśe avasthitavān}.

\begin{cousinweb}
\textbf{\sk{अवतिष्ठते}} (\emph{avatiṣṭhate}) is built from \emph{ava}, \textquotedblleft{}down\textquotedblright{}, and the root √sthā, \textquotedblleft{}to stand\textquotedblright{}, which also gives \emph{tiṣṭhati}, \textquotedblleft{}stands\textquotedblright{}; MW gives \emph{ava-sthā} as to remain or stay, with the participle \emph{avasthita}, and Pāṇini (1.3.22) gives √sthā the middle endings after \emph{ava} (MW).
\end{cousinweb}

\begin{grammarlens}[title={What you've built}]
One more word for getting around.
\end{grammarlens}

\subsection*{Your turn}
\begin{itemize}
\raggedright
  \item \emph{Say it:} \emph{avatiṣṭhate}
  \item \emph{Say it:} \emph{avatiṣṭhate}, once more
  \item \emph{Say it:} \emph{saḥ videśe ekaṃ saṃvatsaram avatiṣṭhate.} — He stays on abroad for a year.
  \item \emph{Recall:} say the Sanskrit for crowded, then the Sanskrit for a customs house, then say \emph{avatiṣṭhate} again
\end{itemize}

\begin{tcolorbox}[breakable,colback=teal!4,colframe=teal!35!black,title={Before you move on}]
What does \sk{अवतिष्ठते} mean? (\textquotedblleft{}Stays on; remains\textquotedblright{}.) Say it once more.
\end{tcolorbox}
\section[\texorpdfstring{\sk{प्रतिश्रयः} (pratiśrayaḥ)}{प्रतिश्रयः (pratiśrayaḥ)}]{\sk{प्रतिश्रयः} (pratiśrayaḥ) — a shelter; a place of refuge}

\label{lesson:SA-C426-pratisrayah}



\begin{quote}
Before the new one: say the Sanskrit for a customs house, then the Sanskrit for stays on.
\end{quote}

\subsection*{You'll want to know: \sk{प्रतिश्रयः}}
\textbf{\sk{प्रतिश्रयः}} — \emph{pratiśrayaḥ} — \textquotedblleft{}a shelter; a place of refuge\textquotedblright{}.

\begin{quote}
\sk{विदेशे} \sk{मित्रस्य} \sk{गृहं} \sk{मम} \sk{प्रतिश्रयः} \sk{अस्ति।}
\end{quote}

\begin{quote}
\emph{videśe mitrasya gṛhaṃ mama pratiśrayaḥ asti.}
\end{quote}

\begin{quote}
Abroad, my friend's house is my shelter.
\end{quote}

A masculine noun.

\begin{cousinweb}
\textbf{\sk{प्रतिश्रयः}} (\emph{pratiśrayaḥ}) is built from \emph{prati}, \textquotedblleft{}towards\textquotedblright{}, and the root √śri, \textquotedblleft{}to lean on, to resort to\textquotedblright{}; MW gives \emph{pratiśraya} as refuge, a place of shelter, a house (MW).
\end{cousinweb}

\begin{grammarlens}[title={What you've built}]
One more word for getting around.
\end{grammarlens}

\subsection*{Your turn}
\begin{itemize}
\raggedright
  \item \emph{Say it:} \emph{pratiśrayaḥ}
  \item \emph{Say it:} \emph{pratiśrayaḥ}, once more
  \item \emph{Say it:} \emph{videśe mitrasya gṛhaṃ mama pratiśrayaḥ asti.} — Abroad, my friend's house is my shelter.
  \item \emph{Recall:} say the Sanskrit for a customs house, then the Sanskrit for stays on, then say \emph{pratiśrayaḥ} again
\end{itemize}

\begin{tcolorbox}[breakable,colback=teal!4,colframe=teal!35!black,title={Before you move on}]
What does \sk{प्रतिश्रयः} mean? (\textquotedblleft{}A shelter; a place of refuge\textquotedblright{}.) Say it once more.
\end{tcolorbox}
